\section{Inverse-incidence coarea on the original Lorentz source}
\label{sec:v61-inverse-incidence}
We return to the physical boundary sources. The contact determinant
contains a gain which is lost in the unweighted finite-count height
bound: it cancels one inverse incidence at every contact, simultaneously.
This gives a density estimate for an unbounded observable of the actual
orbit. The estimate is uniform in the table and in the incidence
threshold, but its constant is exponential in the collision count.
We keep that dependence explicit.

Throughout this section the table is circular, $R\in[9/20,47/100]$.
Write $\ell_0=3/50$, $c_*=91/(10000\pi)$ and
$c_j(x)=\cos\varphi(T_R^jx)$. On a regular $m$-flight orbit put
\begin{equation}\label{eq:v61-inverse-incidence-weight}
 \mathcal I_{m,R}(x)=\prod_{j=0}^{m}c_j(x)^{-1},\qquad
 L_{m,R}(x)=\sum_{j=0}^{m-1}\tau_R(T_R^jx).
\end{equation}
Values on singular orbits may be set to zero. These orbits form a
collision-null set at each finite count. There is no added random
variable or modified flight in \eqref{eq:v61-inverse-incidence-weight}.

\subsection{Cancellation in the endpoint density}
For a fixed regular word, use oriented endpoint arclengths $(u,v)$
and its reduced action $F(u,v)$. Put $t_j=1/\ell_j$. The internal
contact matrix is
\[
 \begin{aligned}
 A_{jj}&=t_{j-1}+t_j+\frac{2}{Rc_j} &&(1\le j<m),\\
 A_{j,j+1}=A_{j+1,j}&=-t_j &&(1\le j<m-1).
 \end{aligned}
\]
An empty determinant, at $m=1$, is one. This is the matrix in
\eqref{eq:v45-contact-matrix}, evaluated at the original contacts.

\begin{lemma}[A determinant gain at every incident contact]
\label{lem:v61-determinant-gain}
On the full regular endpoint image,
\begin{equation}\label{eq:v61-cross-and-product}
 |F_{uv}|=\frac{c_0c_m\prod_{j=0}^{m-1}t_j}{\det A},\qquad
 \mathcal I_{m,R}|F_{uv}|
       \le \left(\frac R2\right)^{m-1}\prod_{j=0}^{m-1}t_j.
\end{equation}
For $S\subset\{1,\ldots,m-1\}$, let $A^{(S)}$ denote the same
matrix with $c_j$ replaced by one for $j\in S$, solely in its
indicated diagonal potentials. Then
\begin{equation}\label{eq:v61-selected-gain}
 \frac{\det A^{(S)}}{\det A}
       \le\prod_{j\in S}\frac{53c_j}{6+47c_j}.
\end{equation}
The matrix $A^{(S)}$ is an algebraic comparison, not the Hessian of
an asserted replacement trajectory.
\end{lemma}
\begin{proof}
The first variation gives $p_0=-F_u$. Hence the absolute Jacobian
of $(u,p_0)\mapsto(u,v)$ is $|F_{uv}|$ in the inverse direction.
The off-diagonal entries of the full contact Hessian have magnitudes
$c_jc_{j+1}t_j$. Its internal block is $D_cAD_c$, where
$D_c=\operatorname{diag}(c_1,\ldots,c_{m-1})$. Taking the endpoint
Schur complement and using the corner cofactor of a tridiagonal
matrix gives the first identity in \eqref{eq:v61-cross-and-product}.
The factors $\prod_{j=1}^{m-1}c_j^2$ cancel. This is the identity
\eqref{eq:v45-cross-density}; its algebra needs positive incidence,
not a protected collar or a critical endpoint. The one-flight case
is the direct mixed derivative of the segment length.

Write $A=L+D$, where $D_{jj}=2/(Rc_j)$ and $L$ is the weighted
Dirichlet path Laplacian, with diagonal $t_{j-1}+t_j$ and neighboring
entries $-t_j$. Its quadratic form is a sum of nonnegative squares,
including the two boundary terms. Thus $L\ge0$, and
\[
 \det A=\det D\,\det(I+D^{-1/2}LD^{-1/2})
       \ge\prod_{j=1}^{m-1}\frac2{Rc_j}.
\]
Substitution proves the second identity's inequality, including
both endpoint incidence factors.

For the refined assertion change one selected diagonal at a time.
Before restoring $c_j$ from one to its physical value, its Schur
complement $s_j$ is positive and no larger than
$t_{j-1}+t_j+2/R\le2/\ell_0+2/R$. Restoring it adds
$(2/R)(c_j^{-1}-1)$ to this complement. The corresponding inverse
determinant ratio is at most
\[
 \frac{(1+R/\ell_0)c_j}{1+(R/\ell_0)c_j}
       \le \frac{53c_j}{6+47c_j}.
\]
The last expression uses $R/\ell_0\le47/6$; the ratio is increasing
in that parameter for $0<c_j\le1$. All intermediate matrices remain
positive definite. Multiplication of the ratios proves
\eqref{eq:v61-selected-gain}, without a factor depending on the
order or number of selected contacts.
\end{proof}

\subsection{An unbounded physical insertion at fixed roof}
\begin{lemma}[Endpoint level integrals]
\label{lem:v61-endpoint-level-integral}
On each of a fixed finite collection of rational endpoint charts,
let $f$ be the reduced action of a regular $m$-flight word. There
are $C>0$ and $A_0>1$ such that, for almost every $t$,
\begin{equation}\label{eq:v61-level-integral}
 \int_{\{f=t\}}\frac{d\mathcal H^1}{|\nabla f|}\le CA_0^m.
\end{equation}
The integral includes the entire open physical endpoint image in the
chart. In particular it does not require that image to be convex.
\end{lemma}
\begin{proof}
Take charts $\alpha=\alpha_*+2\arctan x$, $|x|<1$, at both
endpoints and use the disjoint refinements of
Lemma~\ref{lem:v44-level-complexity}. Their derivatives lie in
$[1,2]$. The first variation is
$\nabla f=R(-a(x)p_0,a(y)p_m)$. On the region where either
$|p_0|$ or $|p_m|$ is at least $\delta_0=1/20$, its magnitude
is at least $R_-\delta_0$.

The level-length argument of that lemma also applies to these endpoint
charts: impose an endpoint coordinate line, rather than an initial
coordinate line, on the same auxiliary collision graph. It still has
$O(m+1)$ variables and bounded-degree polynomial sign conditions.
The coefficient-independent bound \eqref{eq:v44-sign-bound} bounds
each finite projection fiber by $CA_0^m$. Integrating coordinate
multiplicities bounds the total endpoint level length by $CA_0^m$.
Thus the asserted integral over the non-normal region has this bound.

On the complementary region the rational-coordinate Hessian satisfies
$\nabla^2f\ge\lambda_0 I$ with $\lambda_0=R_-/2$, by
\eqref{eq:v44-chart-convexity}. At a regular level its signed
curvature with normal $\nabla f/|\nabla f|$ is
$k_f=T^{\mathsf T}\nabla^2fT/|\nabla f|$, and consequently
$|\nabla f|^{-1}\le |k_f|/\lambda_0$. The total absolute
curvature bound of Lemma~\ref{lem:v44-level-complexity} gives
\eqref{eq:v61-level-integral} on this part. It includes every
connected portion, even if the level meets a physical seam in its
closure. Compact regular exhaustions and monotone convergence retain
all such portions. Critical roof values and exceptional coarea levels
form a null set at each word; no positive-measure roof set is removed.
\end{proof}

\begin{theorem}[Complete inverse-incidence roof density]
\label{thm:v61-inverse-incidence-density}
The measure $\mathcal I_{m,R}\nu$ is finite. Its pushforward by
$L_{m,R}$ has a density $\rho^{\mathcal I}_{m,R}$ satisfying
\begin{equation}\label{eq:v61-inverse-incidence-height}
 \|\rho^{\mathcal I}_{m,R}\|_\infty\le CA^m
       \qquad(m\ge1,\ R\in[9/20,47/100])
\end{equation}
for fixed $C>0$, $A>1$. The same bound holds for every nonnegative
weight $\prod_{j=0}^m c_j^{-\beta_j}$ with $0\le\beta_j\le1$.
For complex insertions $w_{n,R}$ of supremum at most $M$, the exact
actual-return coefficients obey
\begin{equation}\label{eq:v61-exact-weighted-incidence}
 \operatorname*{ess\,sup}_t
 \sum_{n=1}^{m}\sum_{k\in\Z^2}
 |p^{\,w\mathcal I}_{n,R}(k,m,t)|\le\frac{CM A^m}{c_*}.
\end{equation}
Here the insertion on $\{N_{n,R}=m\}$ is the product along that
very $m$-collision orbit. Occupation is counted at $0,\ldots,m-1$,
and terminal membership at $m$ is retained separately.
\end{theorem}
\begin{proof}
Under the full collision probability the endpoint density is
$|F_{uv}|/(4\pi R)$. On a rational endpoint chart the further
Jacobian is $R^2a(x)a(y)$, bounded uniformly above. By
Lemma~\ref{lem:v61-determinant-gain}, multiplication by
$\mathcal I_{m,R}$ leaves a chart density bounded by a constant
times
\[
 \ell_0^{-1}\max\{1,R_+/(2\ell_0)\}^{m-1}.
\]
The bound holds on the full regular word, without a lower cutoff
on any $c_j$.

The finite horizon gives a fixed finite alphabet of next centers.
There are at most $D^m$ words and a fixed number of endpoint charts
per word. Endpoint injectivity is
Lemma~\ref{lem:v44-endpoint-curvature}. The bounded chart density
therefore first proves finiteness of the weighted source, by integration
over the bounded endpoint boxes and summation over words. Coarea on
those same boxes, followed by Lemma~\ref{lem:v61-endpoint-level-integral},
gives \eqref{eq:v61-inverse-incidence-height}. Alternatively one may
truncate the insertion and pass monotonically to its full value; every
bound is independent of the truncation.

Since $c_j\le1$, the fractional-power products are dominated by
$\mathcal I_{m,R}$. For fixed $m$ the original section events
$\Pi_{n,k,m,R}$ are pairwise disjoint as $n,k$ vary. Their laws
are restrictions of $\nu/c_*$, and on each such event $T_{n,R}=L_{m,R}$.
Positive domination before pushforward proves
\eqref{eq:v61-exact-weighted-incidence}, including arbitrary bounded
measurable $w_{n,R}$. This is a sum of variation densities, not a
probability total-variation distance; no factor $1/2$ occurs.
\end{proof}

\begin{corollary}[The original first-incidence source has a height gain]
\label{cor:v61-incidence-height-gain}
Let $b^{\varepsilon,\mathrm{inc},w}_{n,k,m,R}$ be exactly the
first-physical-defect incidence source in
\eqref{eq:v48-two-physical-types}. For all sufficiently small
$\varepsilon>0$, simultaneously for all finite $m\ge2$,
\begin{equation}\label{eq:v61-incidence-height-gain}
 \operatorname*{ess\,sup}_t\sum_{n=1}^{m}\sum_{k\in\Z^2}
 |b^{\varepsilon,\mathrm{inc},w}_{n,k,m,R}(t)|
                 \le C M A^m\varepsilon.
\end{equation}
More generally, restricting the original source to
$c_j<s_j$ at each contact in a prescribed set $S\subset\{0,\ldots,m\}$
costs a factor $\prod_{j\in S}s_j$ in the same height bound.
All constants include the fixed section normalization once.
\end{corollary}
\begin{proof}
The smooth physical guard is between zero and one. A first incidence
factor below one, including its transition region, has
$c_j<2\varepsilon q_0^{d_j/4}$, where
$q_0=47/53$ and $d_j=\min(j,m-j)$. Pointwise on the unchanged orbit,
\[
 (1-H_{m,R}^{\varepsilon})\one_{\{\text{first type is incidence}\}}
 \le\sum_{j=0}^{m}\one_{\{c_j<2\varepsilon q_0^{d_j/4}\}}
 \le2\varepsilon\sum_{j=0}^{m}q_0^{d_j/4}\mathcal I_{m,R}.
\]
The geometric sum is at most $2/(1-q_0^{1/4})$. Apply
Theorem~\ref{thm:v61-inverse-incidence-density} and source domination.
For simultaneous contact conditions use
$\prod_{j\in S}\one_{\{c_j<s_j\}}
 \le(\prod_{j\in S}s_j)\prod_{j\in S}c_j^{-1}
 \le(\prod_{j\in S}s_j)\mathcal I_{m,R}$.
The first-defect source is only enlarged for a positive upper bound;
its actual return index, collision count and displacement are not changed.
\end{proof}

\paragraph{Collision-count dependence.}
Equation~\eqref{eq:v61-incidence-height-gain} is a density estimate,
not the mass estimate $O(\varepsilon^2)$ of the incidence layer.
It allows an arbitrary count-dependent $\varepsilon_m$. In particular,
its normalized height tends to zero if
$m^2A^m\varepsilon_m\to0$. For a central exact label whose canonical
reference $G(u)$ is at least $d>0$, the pointwise lower law
\eqref{eq:v51-lower-local-law} gives $P(u)\ge d/2$ for large $m$.
Division at the same roof then bounds the original conditional incidence
source by $C_d m^2A^m\varepsilon$ almost everywhere there. A single
source disintegration supplies these inequalities for all bounded tests.

This exponential factor is not a central-scale height estimate for a
fixed $\varepsilon$. In particular it does not establish the ordered
limit in Corollary~\ref{cor:v51-positive-criterion}. The missing step
is uniform control of the positive exact-label level sum as $m$ grows,
not integrability of the inverse incidence at a fixed count. No change
of the finite arithmetic coefficient is involved.
